% titel: Blatt 8, A9: CsCl-Kantenlänge und kürzester Cs-Cl-Abstand
% skript: Mortimer Kap. 12 (Elementarzellen, CsCl-Typ)

% Teil I: Werkzeuge, kein Lösungsweg.

Aus der Dichte eines Salzes kann man die Zellgeometrie und Ionenabstände
bestimmen, wenn der Strukturtyp bekannt ist.

\begin{werkzeugkasten}[CsCl-Struktur]
Im CsCl-Typ (`CsCl`, `kubisch-primitiv` mit zweiatomiger Basis) gilt:
\begin{itemize}
  \item \emph{Eine} Formeleinheit pro Elementarzelle ($Z=1$).
  \item Cl$^-$ auf den 8 Ecken, Cs$^+$ in der Wurffmitte (oder umgekehrt).
  \item Kantenlänge $a$; die beiden Ionen liegen auf der Raumdiagonale.
  \item Dichte:
    \[
      \rho = \frac{Z\,M}{N_A\,a^3}\qquad\text{also}\qquad
      a = \left(\frac{Z\,M}{N_A\,\rho}\right)^{1/3}.
    \]
\end{itemize}
Kürzester Ionenabstand = halbe Raumdiagonale $= \frac{a\sqrt{3}}{2}$.
\textit{Skript:} Mortimer Kap. 12 (Elementarzellen).
\end{werkzeugkasten}

\begin{achtung}[Darauf achten]
\begin{itemize}
  \item $M(\mathrm{CsCl}) \approx 168.36\ g\,mol^{-1}$ (Cs 132.9 $+$ Cl 35.45).
  \item Einheiten: $\rho = 3.99\ g\,cm^{-3}$, $a$ erhält man in cm, dann in pm
        ($1\ cm = 10^{10}\ pm$).
  \item Im CsCl-Gitter ist es \emph{die} Raumdiagonale, nicht die Kante, auf
        der sich die Ionen berühren.
\end{itemize}
\end{achtung}
